\documentclass[tikz,border=8pt]{standalone}
\usepackage{ctex,amsmath,amssymb}
\usetikzlibrary{arrows.meta,positioning,calc,fit,backgrounds}
% Shared visual language for LFS architecture diagrams.
\RequirePackage{../../mymath}

\definecolor{LFSBlue}{HTML}{2563EB}
\definecolor{LFSBlueSoft}{HTML}{DBEAFE}
\definecolor{LFSGreen}{HTML}{15803D}
\definecolor{LFSGreenSoft}{HTML}{DCFCE7}
\definecolor{LFSOrange}{HTML}{C2410C}
\definecolor{LFSOrangeSoft}{HTML}{FFEDD5}
\definecolor{LFSRed}{HTML}{B91C1C}
\definecolor{LFSRedSoft}{HTML}{FEE2E2}
\definecolor{LFSPurple}{HTML}{7E22CE}
\definecolor{LFSPurpleSoft}{HTML}{F3E8FF}
\definecolor{LFSGray}{HTML}{475569}
\definecolor{LFSGraySoft}{HTML}{F1F5F9}

\tikzset{
  lfs picture/.style={
    font=\small,
    node distance=8mm and 12mm,
    >=Latex,
    every path/.style={line cap=round, line join=round}
  },
  block/.style={
    draw=LFSBlue,
    fill=LFSBlueSoft,
    rounded corners=2pt,
    line width=0.8pt,
    align=center,
    inner xsep=7pt,
    inner ysep=5pt,
    minimum height=10mm
  },
  compute/.style={
    block,
    draw=LFSPurple,
    fill=LFSPurpleSoft
  },
  storage/.style={
    block,
    draw=LFSGreen,
    fill=LFSGreenSoft
  },
  interface/.style={
    block,
    draw=LFSOrange,
    fill=LFSOrangeSoft
  },
  risk/.style={
    block,
    draw=LFSRed,
    fill=LFSRedSoft
  },
  neutral/.style={
    block,
    draw=LFSGray,
    fill=LFSGraySoft
  },
  decision/.style={
    diamond,
    draw=LFSOrange,
    fill=LFSOrangeSoft,
    line width=0.8pt,
    align=center,
    inner sep=2pt,
    aspect=2
  },
  group/.style={
    draw=LFSGray,
    rounded corners=3pt,
    dashed,
    line width=0.7pt,
    inner sep=6pt
  },
  arrow/.style={
    -{Latex[length=2.4mm,width=1.6mm]},
    draw=LFSGray,
    line width=0.9pt
  },
  data arrow/.style={
    arrow,
    draw=LFSBlue
  },
  control arrow/.style={
    arrow,
    draw=LFSOrange,
    dashed
  },
  residual/.style={
    arrow,
    draw=LFSGreen,
    rounded corners=4pt
  },
  label text/.style={
    font=\scriptsize,
    text=LFSGray,
    fill=white,
    inner sep=1.5pt
  },
  title/.style={
    font=\bfseries,
    text=LFSGray,
    align=center
  }
}


\begin{document}
\begin{tikzpicture}[my picture, node distance=5mm and 6.5mm]

% 顶部：资源预算与算力基础设施
\node[neutral, text width=70mm, minimum height=13mm] (hw_infra) {
  \textbf{计算基础设施与通信拓扑}\\[0.8mm]
  \scriptsize AI 加速卡集群架构、高带宽显存（HBM）\\
  NVLink 节点内互联与 InfiniBand 跨节点通信（第61章）
};
\node[storage, text width=70mm, minimum height=13mm, right=5mm of hw_infra] (mem_budget) {
  \textbf{显存分析与资源预算建模}\\[0.8mm]
  \scriptsize 静态参数、梯度、优化器状态与动态激活显存\\
  训练与推理吞吐、通信开销理论建模（第32章）
};
\begin{scope}[on background layer]
  \node[group, fit=(hw_infra)(mem_budget), label={[font=\footnotesize\bfseries, text=MyGray]above:一、硬件基础设施与显存建模底座}] (g_base) {};
\end{scope}

% 中部：分布式训练系统
\node[compute, text width=70mm, minimum height=13mm, below=11mm of hw_infra.south west, anchor=north west] (train_3d) {
  \textbf{3D 混合分布式并行拓扑}\\[0.8mm]
  \scriptsize 数据并行（DP）、张量并行（TP）与流水线并行（PP）\\
  序列并行与通信计算重叠隐藏（第62章）
};
\node[compute, text width=70mm, minimum height=13mm, right=5mm of train_3d] (train_zero) {
  \textbf{ZeRO 显存分片与大规模训练}\\[0.8mm]
  \scriptsize 优化器、梯度与参数多级切分（ZeRO-1/2/3）\\
  长序列通信与容灾容错架构（第62章）
};
\begin{scope}[on background layer]
  \node[group, fit=(train_3d)(train_zero), label={[font=\footnotesize\bfseries, text=MyPurple]above:二、大规模分布式训练系统}] (g_train) {};
\end{scope}

% 底部：推理服务与加速运行时
\node[interface, text width=45mm, minimum height=14mm, below=11mm of train_3d.south west, anchor=north west] (inf_serv) {
  \textbf{推理服务与动态调度}\\[0.8mm]
  \scriptsize 连续批处理（Continuous Batching）\\
  Prefill/Decode 算力解耦（第63、64章）
};
\node[storage, text width=45mm, minimum height=14mm, right=5mm of inf_serv] (inf_cache) {
  \textbf{增量解码与分页缓存}\\[0.8mm]
  \scriptsize PagedAttention 显存虚拟化\\
  KV Cache 显存零碎片调度（第51章）
};
\node[block, text width=45mm, minimum height=14mm, right=5mm of inf_cache] (inf_opt) {
  \textbf{算法加速与模型量化}\\[0.8mm]
  \scriptsize 投机解码与算子融合\\
  AWQ/GPTQ/FP8 低比特量化（第52、53章）
};
\begin{scope}[on background layer]
  \node[group, fit=(inf_serv)(inf_cache)(inf_opt), label={[font=\footnotesize\bfseries, text=MyOrange]above:三、高性能推理服务运行时与算子优化}] (g_inf) {};
\end{scope}

% 层级流动连线
\draw[data arrow] ($(hw_infra.south)+(15mm,0)$) -- ($(train_3d.north)+(15mm,0)$);
\draw[data arrow] ($(mem_budget.south)-(15mm,0)$) -- ($(train_zero.north)-(15mm,0)$);

\draw[data arrow] (train_3d.east) -- (train_zero.west);
\draw[data arrow] (inf_serv.east) -- (inf_cache.west);
\draw[data arrow] (inf_cache.east) -- (inf_opt.west);

\draw[data arrow] ($(train_3d.south)-(10mm,0)$) -- ($(inf_serv.north)+(0,0)$);
\draw[data arrow] ($(train_zero.south)+(10mm,0)$) -- ($(inf_opt.north)+(0,0)$);

% 底部总结栏
\node[neutral, text width=146mm, minimum height=7mm, below=6mm of inf_cache.south] (bot_summary) {
  \scriptsize \textbf{第四篇系统主线：} 芯片与网络基建 $\longrightarrow$ 显存与通信预算 $\longrightarrow$ 3D分布式并行训练 $\longrightarrow$ 动态批处理服务 $\longrightarrow$ KV缓存虚拟化与量化加速
};

\end{tikzpicture}
\end{document}
